\documentclass[11pt]{article}
\newcommand{\ListaTitulo}{Lista 12: Distribuições Uniforme e Normal}
% Preâmbulo compartilhado: MATF14, Listas de Exercícios 2026
% Usado por ListaNN.tex e GabaritoNN.tex via % Preâmbulo compartilhado: MATF14, Listas de Exercícios 2026
% Usado por ListaNN.tex e GabaritoNN.tex via % Preâmbulo compartilhado: MATF14, Listas de Exercícios 2026
% Usado por ListaNN.tex e GabaritoNN.tex via \input{preamble}

\usepackage[utf8]{inputenc}
\usepackage[T1]{fontenc}
\usepackage[brazil]{babel}
\usepackage{fullpage}
\usepackage{amsmath,amssymb}
\usepackage{enumitem}
\usepackage{xcolor}
\usepackage{fancyhdr}
\usepackage{titlesec}
\usepackage{listings}
\usepackage{hyperref}
\usepackage{booktabs}
\usepackage{graphicx}
\usepackage{tikz}

\definecolor{matf14azul}{HTML}{0D3B54}
\definecolor{matf14dourado}{HTML}{C9971E}

\titleformat{\section}{\normalfont\Large\bfseries\color{matf14azul}}{\thesection}{1em}{}
\titleformat{\subsection}{\normalfont\large\bfseries\color{matf14azul}}{\thesubsection}{1em}{}

\providecommand{\ListaTitulo}{Lista de Exercícios}

\setlength{\headheight}{14pt}
\pagestyle{fancy}
\fancyhf{}
\lhead{\small MATF14: Estatística Econômica I}
\rhead{\small \ListaTitulo}
\cfoot{\thepage}
\renewcommand{\headrulewidth}{0.4pt}

\lstset{
  basicstyle=\ttfamily\small,
  breaklines=true,
  frame=single,
  columns=fullflexible,
  backgroundcolor=\color{gray!8},
  commentstyle=\color{matf14dourado},
  keywordstyle=\color{matf14azul}\bfseries,
  language=R
}

\newcounter{questaocounter}
\newenvironment{questao}[1]{%
  \stepcounter{questaocounter}%
  \bigskip\noindent\textbf{Questão #1.}\ %
}{\par}

\newcommand{\cabecalho}[2]{%
  \begin{center}
    {\Large\bfseries\color{matf14azul} #1}\\[0.3em]
    {\large #2}
  \end{center}
  \vspace{0.5em}
  \noindent\rule{\linewidth}{0.6pt}
  \vspace{0.5em}
}

\newenvironment{solucao}{%
  \par\smallskip\noindent\textit{\color{matf14dourado!80!black}Solução.}\ %
}{\hfill$\blacksquare$\par}


\usepackage[utf8]{inputenc}
\usepackage[T1]{fontenc}
\usepackage[brazil]{babel}
\usepackage{fullpage}
\usepackage{amsmath,amssymb}
\usepackage{enumitem}
\usepackage{xcolor}
\usepackage{fancyhdr}
\usepackage{titlesec}
\usepackage{listings}
\usepackage{hyperref}
\usepackage{booktabs}
\usepackage{graphicx}
\usepackage{tikz}

\definecolor{matf14azul}{HTML}{0D3B54}
\definecolor{matf14dourado}{HTML}{C9971E}

\titleformat{\section}{\normalfont\Large\bfseries\color{matf14azul}}{\thesection}{1em}{}
\titleformat{\subsection}{\normalfont\large\bfseries\color{matf14azul}}{\thesubsection}{1em}{}

\providecommand{\ListaTitulo}{Lista de Exercícios}

\setlength{\headheight}{14pt}
\pagestyle{fancy}
\fancyhf{}
\lhead{\small MATF14: Estatística Econômica I}
\rhead{\small \ListaTitulo}
\cfoot{\thepage}
\renewcommand{\headrulewidth}{0.4pt}

\lstset{
  basicstyle=\ttfamily\small,
  breaklines=true,
  frame=single,
  columns=fullflexible,
  backgroundcolor=\color{gray!8},
  commentstyle=\color{matf14dourado},
  keywordstyle=\color{matf14azul}\bfseries,
  language=R
}

\newcounter{questaocounter}
\newenvironment{questao}[1]{%
  \stepcounter{questaocounter}%
  \bigskip\noindent\textbf{Questão #1.}\ %
}{\par}

\newcommand{\cabecalho}[2]{%
  \begin{center}
    {\Large\bfseries\color{matf14azul} #1}\\[0.3em]
    {\large #2}
  \end{center}
  \vspace{0.5em}
  \noindent\rule{\linewidth}{0.6pt}
  \vspace{0.5em}
}

\newenvironment{solucao}{%
  \par\smallskip\noindent\textit{\color{matf14dourado!80!black}Solução.}\ %
}{\hfill$\blacksquare$\par}


\usepackage[utf8]{inputenc}
\usepackage[T1]{fontenc}
\usepackage[brazil]{babel}
\usepackage{fullpage}
\usepackage{amsmath,amssymb}
\usepackage{enumitem}
\usepackage{xcolor}
\usepackage{fancyhdr}
\usepackage{titlesec}
\usepackage{listings}
\usepackage{hyperref}
\usepackage{booktabs}
\usepackage{graphicx}
\usepackage{tikz}

\definecolor{matf14azul}{HTML}{0D3B54}
\definecolor{matf14dourado}{HTML}{C9971E}

\titleformat{\section}{\normalfont\Large\bfseries\color{matf14azul}}{\thesection}{1em}{}
\titleformat{\subsection}{\normalfont\large\bfseries\color{matf14azul}}{\thesubsection}{1em}{}

\providecommand{\ListaTitulo}{Lista de Exercícios}

\setlength{\headheight}{14pt}
\pagestyle{fancy}
\fancyhf{}
\lhead{\small MATF14: Estatística Econômica I}
\rhead{\small \ListaTitulo}
\cfoot{\thepage}
\renewcommand{\headrulewidth}{0.4pt}

\lstset{
  basicstyle=\ttfamily\small,
  breaklines=true,
  frame=single,
  columns=fullflexible,
  backgroundcolor=\color{gray!8},
  commentstyle=\color{matf14dourado},
  keywordstyle=\color{matf14azul}\bfseries,
  language=R
}

\newcounter{questaocounter}
\newenvironment{questao}[1]{%
  \stepcounter{questaocounter}%
  \bigskip\noindent\textbf{Questão #1.}\ %
}{\par}

\newcommand{\cabecalho}[2]{%
  \begin{center}
    {\Large\bfseries\color{matf14azul} #1}\\[0.3em]
    {\large #2}
  \end{center}
  \vspace{0.5em}
  \noindent\rule{\linewidth}{0.6pt}
  \vspace{0.5em}
}

\newenvironment{solucao}{%
  \par\smallskip\noindent\textit{\color{matf14dourado!80!black}Solução.}\ %
}{\hfill$\blacksquare$\par}


\begin{document}

\cabecalho{Lista de Exercícios 12}{Distribuições Uniforme e Normal (Cap. 4, Aulas 27--28)}

\noindent \textit{Justifique todas as respostas. As resoluções são discutidas em sala e nos horários de atendimento; não são publicadas neste material.}

\begin{questao}{1} \textbf{(Uniforme: horário de chegada de um ônibus)}

Um ônibus passa em um ponto a cada 15 minutos, sem horário fixo divulgado. O tempo de espera $X$
de um passageiro que chega ao ponto em um instante aleatório é $X\sim\text{Unif}(0,15)$.

\begin{enumerate}[label=(\alph*)]
  \item Calcule $E(X)$ e interprete: por que faz sentido que o tempo médio de espera seja metade
    do intervalo entre ônibus?
  \item Calcule $P(X>10)$.
  \item Calcule $dp(X)$.
\end{enumerate}
\end{questao}

\begin{questao}{2} \textbf{(Normal: notas de uma prova)}

As notas de uma prova de estatística em um curso de 200 alunos seguem aproximadamente
$X\sim N(6{,}5,\ 1{,}2^2)$.

\begin{enumerate}[label=(\alph*)]
  \item Qual a proporção esperada de alunos com nota abaixo de 5,0 (nota mínima para aprovação)?
  \item O professor quer dar menção "Excelência" aos 10\% com nota mais alta. Qual nota mínima
    corresponde a esse corte?
  \item Estime quantos alunos, dos 200, devem ficar abaixo da média da turma. Esse número deveria
    surpreender alguém que já conhece a propriedade de simetria da Normal? Justifique.
\end{enumerate}
\end{questao}

\begin{questao}{3} \textbf{(Padronização: comparando desempenho em escalas diferentes)}

Um candidato obteve 720 pontos em uma prova de matemática, cuja distribuição na população de
candidatos é $N(650, 80^2)$, e 78 pontos em uma prova de redação, com distribuição $N(70, 10^2)$.

\begin{enumerate}[label=(\alph*)]
  \item Calcule o escore padronizado ($z$) do candidato em cada prova.
  \item Em qual das duas provas o candidato teve um desempenho relativamente melhor frente aos
    demais candidatos? Justifique usando os escores padronizados, não as notas brutas.
  \item Por que comparar diretamente 720 e 78 (as notas brutas) seria uma comparação sem sentido?
\end{enumerate}
\end{questao}

\begin{questao}{4} \textbf{(Regra empírica: controle de qualidade de peso)}

O peso de pacotes de um produto alimentício é $X\sim N(500\text{g},\ 5^2)$ (rótulo indica 500g).

\begin{enumerate}[label=(\alph*)]
  \item Usando a regra empírica (68-95-99,7), em que intervalo de peso está aproximadamente 95\%
    dos pacotes, sem usar `pnorm`/tabela, só a regra?
  \item Calcule $P(490 < X < 510)$ exatamente (via padronização) e compare com o valor
    aproximado do item (a).
  \item Um órgão de defesa do consumidor reclama de pacotes "muito abaixo do rótulo". Calcule
    $P(X<485)$ e avalie se essa reclamação parece estatisticamente justificada, dado o modelo.
\end{enumerate}
\end{questao}

\begin{questao}{5} \textbf{(Dedução: esperança e variância da Uniforme)}

Seja $X\sim\text{Unif}(a,b)$, com densidade $f(x)=\frac{1}{b-a}$ para $a\le x\le b$.

\begin{enumerate}[label=(\alph*)]
  \item Mostre, por integração direta, que $E(X) = \dfrac{a+b}{2}$.
  \item Mostre, por integração direta, que $\mathrm{Var}(X) = \dfrac{(b-a)^2}{12}$. (Dica: use a
    fórmula alternativa $\mathrm{Var}(X)=E(X^2)-[E(X)]^2$, calculando primeiro
    $E(X^2)=\int_a^b x^2\cdot\frac{1}{b-a}dx$.)
\end{enumerate}
\end{questao}

\bigskip
\noindent\textit{A questão a seguir não tem resposta única e não é cobrada em prova no mesmo
formato: fecha a Unidade 4 e é para discussão em grupo ou por escrito.}

\begin{questao}{6 (discussão: síntese do curso)} \textbf{(Sem resposta única)}

Retornos financeiros reais costumam ter curtose maior que 3 (caudas mais pesadas que a Normal,
Capítulo 1). Escreva um parágrafo explicando, para um investidor não técnico, por que usar a
distribuição Normal para estimar a probabilidade de uma perda extrema pode ser perigoso, e o que
"caudas gordas" significam em termos práticos de risco.
\end{questao}


\bigskip
\section*{Exercícios complementares}

\begin{enumerate}[label=\textbf{C\arabic*.}, leftmargin=*]
\item Sejam $a$ e $\theta$ positivos e $g(y) = 1/\theta$ para $a \leq y \leq a + \theta$, zero fora.
\begin{enumerate}[label=(\alph*)]
  \item Identifique a distribuição de $Y$.
  \item Calcule $E(Y)$ e $Var(Y)$.
  \item $Var(Y)$ depende de $\theta$ mas não de $a$. Interprete.
\end{enumerate}
\item A mediana de uma variável contínua com FDA $F$ é o valor $m$ com $F(m) = 1/2$. Encontre a mediana de $X \sim Unif(\alpha, \beta)$ e de $X \sim N(\mu, \sigma^2)$.
\item As vendas mensais de um produto são normais com média 500 e desvio padrão 50 unidades. Se a empresa fabricar 600 unidades num mês, qual a probabilidade de não atender todos os pedidos?
\item As notas do SAT (exame semelhante ao ENEM nos EUA) são normais com média 500 e desvio padrão 100. Que porcentagem das notas está
\begin{enumerate}[label=(\alph*)]
  \item entre 500 e 600?
  \item entre 400 e 600?
  \item acima de 600?
  \item abaixo de 300?
\end{enumerate}
\item A carga de trabalho de agentes de liberdade condicional é normal com média 115 casos e desvio padrão 10. Qual a probabilidade de
\begin{enumerate}[label=(\alph*)]
  \item um agente ter entre 90 e 105 casos?
  \item um agente ter mais de 135 casos?
  \item quatro agentes independentes terem, todos, mais de 135 casos?
\end{enumerate}
\item O retorno anual de um fundo é normal com média 8\% e desvio padrão 20\%. Calcule a probabilidade de o retorno ser maior que 10\% e de ser menor que $-5\%$.

\end{enumerate}

\bigskip
\needspace{12\baselineskip}
\section*{Aprofundamento: modelos probabilísticos e dados reais}
\noindent\textit{Exercícios ligados à seção de aplicação macroeconômica do Capítulo 4 do livro.}

\begin{enumerate}[label=\textbf{A\arabic*.}, leftmargin=*]
\item \textbf{(Erro de arredondamento e a Uniforme.)} O IBGE divulga as taxas de crescimento do PIB com uma casa decimal. Suponha que o erro de
arredondamento $U$ (valor verdadeiro menos valor divulgado, em pontos percentuais) seja uniforme em $[-0{,}05;\ 0{,}05]$.
\begin{enumerate}[label=(\alph*)]
  \item Calcule $E(U)$, $\mathrm{Var}(U)$ e o desvio padrão de $U$.
  \item Calcule $P(|U| > 0{,}03)$.
  \item A taxa divulgada foi 0,1\%. Qual a probabilidade de a taxa verdadeira ter sido maior que 0,12\%?
\end{enumerate}

\item \textbf{(O crescimento do PIB é normal?)} As 121 taxas de variação do PIB em relação ao trimestre anterior (2º trimestre de 1996 ao 2º trimestre de 2026;
IBGE) têm média $0{,}58\%$ e desvio padrão $1{,}58\%$. O PIB caiu em 27 trimestres; caiu mais de 2\% em 4 trimestres; e 102 taxas ficaram a menos de um desvio padrão da média.
\begin{enumerate}[label=(\alph*)]
  \item Supondo $X \sim N(0{,}58;\ 1{,}58^2)$, calcule $P(X < 0)$, $P(X < -2)$ e $P(|X - 0{,}58| < 1{,}58)$.
  \item Compare com as proporções observadas. Os dados são mais ou menos concentrados em torno da média do que a normal prevê?
  \item A curtose das taxas é 16,0 (a da normal é 3). Relacione esse número com a sua resposta a (b) e com os trimestres de 2020.
\end{enumerate}

\end{enumerate}

\end{document}
